\subsection{Graph of a Relation} 
\begin{definition}[Graph of a Relation]
The \term{graph} of a relation is a picture, in the $xy$ plane, of all the ordered pairs that are elements of the relation.
\end{definition}
\begin{Example}
Sketch the graph of the relation shown below.
\begin{icpic}[x=0.5\linespace,y=0.5\linespace]
\setgraphsquare{5}
\GraphSmall
\regtextright{-5.5}{3}{$R=\left\{(-1,-2),(4,1),(-1,3)\right\}$};
\end{icpic}
\end{Example}
%When a relation is given in set-builder notation, its graph is less obvious.
%\begin{Example}
%The relation $S$ and its graph are shown below.
%\begin{icpic}[x=0.5\linespace,y=0.5\linespace]
%\setgraphsquare{6}
%\GraphSmall
%\draw[graphend] (0,0) circle (5);
%\draw[graphend] (3,4) node[anchor=south west] {$S$};
%\regtextright{-6.5}{4}{$S=\left\{(x,y)\mid x^2+y^2=25\right\}$};
%\end{icpic}
%We can't check that every point on the graph is an ordered pair in $S$ because there are \makeblank of them, but let's check a few.
%\begin{cpic}{}{3}
%\end{cpic}
%\end{Example}
Sometimes, we give a relation as its graph only.
\begin{Example}
The relation $T$ is shown below.
\begin{icpic}[x=0.5\linespace,y=0.5\linespace]
\setgraphsquare{6}
\GraphSmall
\draw[graphend] (-4,-2) -- (-1,4);
\draw[graphend] (-2,1) -- (4,-5);
\filldraw[dotfull] (-4,-2) circle (0.2);
\filldraw[dotempty] (-1,4) circle (0.2);
\filldraw[dotempty] (-2,1) circle (0.2);
\filldraw[dotfull] (4,-5) circle (0.2);
\draw[graphend] (5,5) node[anchor=north east]{$T$};
\end{icpic}
\begin{itemize}
\item The statement $1\relT -2$ is \makeblank[1in] because the point \blankpair\ is on the graph
\item The statement $3\relT 1$ is \makeblank[1in] because the point \blankpair\ is not on the graph
\item The statement $-1\relT 4$ is \makeblank[1in] because the point \blankpair\ is not on the graph of $\relT$ (note that the \makeblank[1in] dot represents a point that \makeblank part of the graph)
\item The statement $4\relT -5$ is \makeblank[1in] because the point \blankpair\ is on the graph of $\relT$ (note that the \makeblank[1in] dot represents a point that \makeblank part of the graph)
\end{itemize}
\end{Example}